We now specialise the apparatus of Section \ref{sec.pointProcesses} to the order flow of
Section \ref{sec.orderDrivenMarkets}.

\begin{defi}\label{def.orderFlowProcess}
 The \emph{order flow} is a marked point process $(\arrivalTimes[]_n, \event[n], \orderSize_n)$
 in which
 \begin{enumerate}[label={\emph{\roman{*}}.}]
  \item $(\arrivalTimes[]_n, \event[n])$ is an exponential-kernel Hawkes process of dimension
        $\numEventTypes = 6$, whose types are the market, limit and withdrawal orders on each
        side, in the order of the table below;
  \item the marks $\orderSize_n$ are the sizes of the orders, independent and identically
        distributed with mean $\meanSize$, independent of the arrival times and of the types.
 \end{enumerate}
\end{defi}

The six types are
\begin{center}
 \begin{tabular}{lcccccc}
  \toprule
  $e$ & market buy & market sell & limit buy & limit sell & withdraw bid & withdraw ask \\
  \midrule
  $\pressure_e$ & $+1$ & $-1$ & $+1$ & $-1$ & $-1$ & $+1$ \\
  \bottomrule
 \end{tabular}
\end{center}
The vector $\pressure$ records the buy/sell pressure exerted on the mid-price:
a market buy consumes the ask and a withdrawal from the bid removes support,
so the first pushes up and the second down.

\begin{defi}\label{def.orderFlowSpecification}
 A \emph{specification} of the order flow is a triple
 $(\baseIntensity, \excitation, \decay)$ of baselines, excitation matrix and decay rate
 for the Hawkes process of Definition \ref{def.orderFlowProcess}.
\end{defi}

\begin{prop}\label{prop.signedSizeContribution}
 Under Assumption \ref{assumption.idealisedBook},
 the contribution \eqref{eq.orderFlowContribution} of the $n$-th event is
 \begin{equation}\label{eq.signedSizeContribution}
  \orderFlowContribution = \pressure_{\event[n]} \, \orderSize_n ,
 \end{equation}
 so that
 $\OFI_{t,\ofiWindow}
 = \sum_{n \, : \, t - \ofiWindow < \arrivalTimes[]_n \leq t}
 \pressure_{\event[n]} \orderSize_n$.
\end{prop}
\begin{proof}
 The first two terms of \eqref{eq.orderFlowContribution} are the signed flow on the bid side
 and the last two are minus the signed flow on the ask side.
 By Assumption \ref{assumption.idealisedBook}~\emph{ii.} no order reaches past the quote it
 acts on, so each event changes one queue by its own size.
 A limit buy adds $\orderSize_n$ to the bid and a withdrawal from the bid removes it,
 giving $\pm\orderSize_n$ with no ask term;
 a market buy removes $\orderSize_n$ from the ask, giving $+\orderSize_n$ through the ask
 terms, and mutatis mutandis for the three sell-side types.
 In each case the sign is $\pressure_{\event[n]}$.
\end{proof}

\begin{defi}\label{def.directionSymmetry}
 Let $\swapMatrix$ be the permutation matrix exchanging the two members of each pair of
 types, so that $\swapMatrix\pressure = -\pressure$.
 A specification is \emph{direction-symmetric} if
 \begin{equation}\label{eq.directionSymmetry}
  \swapMatrix \excitation \swapMatrix = \excitation ,
  \qquad
  \swapMatrix \baseIntensity = \baseIntensity .
 \end{equation}
\end{defi}

\begin{prop}\label{prop.directionSymmetry}
 Let the specification be direction-symmetric.
 Then $\pressure^{\transpose} M v = 0$
 for every matrix $M$ with $\swapMatrix M \swapMatrix = M$
 and every vector $v$ with $\swapMatrix v = v$.
 In particular,
 \begin{equation}\label{eq.baselineSymmetry}
  \pressure^{\transpose}\baseIntensity = 0.
 \end{equation}
 Moreover, if $\branchingRatio < 1$ then
 $\pressure^{\transpose}\stationaryIntensity = 0$.
\end{prop}
\begin{proof}
 $\pressure^{\transpose} M v
 = (\swapMatrix\pressure)^{\transpose}(\swapMatrix M \swapMatrix)(\swapMatrix v)
 \cdot (-1)
 = -\pressure^{\transpose} M v$,
 since $\swapMatrix$ is a symmetric involution and $\swapMatrix\pressure = -\pressure$.
 Taking $M = I$ and $v = \baseIntensity$ gives \eqref{eq.baselineSymmetry}.
 At $\branchingRatio < 1$ the matrix $I - \branchingMatrix$ is invertible and
 $\swapMatrix(I-\branchingMatrix)^{-1}\swapMatrix = (I-\branchingMatrix)^{-1}$,
 so \eqref{eq.stationaryIntensity} gives
 $\pressure^{\transpose}\stationaryIntensity = 0$.
\end{proof}

For a direction-symmetric specification with $\branchingRatio < 1$,
the matrices fixed by $M \mapsto \swapMatrix M \swapMatrix$ are the kernel of a linear map,
hence a closed subspace, and they are closed under products,
so they contain $\excitation$ and $\hurwitzMatrix$ together with the limits and integrals
of their polynomials, among them $e^{\hurwitzMatrix \forecastHorizon}$,
$\meanResponseIntegral(\forecastHorizon)$ and $(I - \branchingMatrix)^{-1}$.
Every expression of Section \ref{sec.forecastingTheMid} built from these and contracted with
$\pressure$ therefore loses its baseline term.

By Proposition \ref{prop.clusterRepresentation},
$\branchingMatrix\subscriptee$ is the expected number of direct offspring of type $e$
born of one event of type $\eone$,
so the column sums $\mathbf{1}^{\transpose}\branchingMatrix$ are readable one entry at a
time.

The kernel carries the mechanism of replenishment.
The two kinds of event that deplete a queue --- an aggression against it and a withdrawal
from it --- must beget successors,
and what they chiefly beget must be the limit orders that refill the side they emptied;
a limit order that has already refilled should beget little.
A kernel that does this has large column sums on the market-order and withdrawal types,
a small one on the limit-order types,
and its mass concentrated off the diagonal blocks that $\pressure$ separates.

When specifying a parametrisation we choose the excitation structure $\branchingMatrix$
and the stationary intensity $\stationaryIntensity$ at which the flow is to run,
and let
\begin{equation}\label{eq.baselineFromTarget}
 \baseIntensity = (I - \branchingMatrix)\stationaryIntensity
\end{equation}
follow.
The pair $(\branchingMatrix, \stationaryIntensity)$ should be chosen with
$\branchingRatio < 1$ and so that the $\baseIntensity$ given by
\eqref{eq.baselineFromTarget} is strictly positive,
these being the hypotheses of Theorem \ref{thm.stationarity}.

Once a parametrisation for the order flow of Definition \ref{def.orderFlowProcess} is chosen,
the intensity informs us of two aspects of the trading epoch:
the direction in which the flow is currently leaning,
and how that direction is assembled from the events that produced it.

\begin{defi}\label{def.intensityContrast}
 Let
 \begin{equation*}
  \upIntensity(t) := \sum_{e \, : \, \pressure_e = +1} \intensity(t),
  \qquad
  \downIntensity(t) := \sum_{e \, : \, \pressure_e = -1} \intensity(t) .
 \end{equation*}
 The \emph{intensity contrast} is
 \begin{equation}\label{eq.intensityContrast}
  \intensityContrast(t) := \upIntensity(t) - \downIntensity(t)
  = \pressure^{\transpose}\intensity[ ](t)
  = \pressure^{\transpose}\baseIntensity
  + \pressure^{\transpose}\excitation \decayedCounts(t) .
 \end{equation}
\end{defi}

Under direction symmetry the first term of \eqref{eq.intensityContrast} vanishes, by
\eqref{eq.baselineSymmetry}, and
\begin{equation}\label{eq.contrastFromState}
 \intensityContrast(t) = \pressure^{\transpose}\excitation \decayedCounts(t) .
\end{equation}
The contrast is then a function of the state alone.

\begin{prop}\label{prop.signedExcitation}
 Let the specification be direction-symmetric, and set
 \begin{equation}\label{eq.signedExcitation}
  \signedExcitation_{\eone}
  := \frac{\big( \pressure^{\transpose}\excitation \big)_{\eone}}{\pressure_{\eone}} ,
  \qquad \eone = 1, \dots, \numEventTypes .
 \end{equation}
 Then $\signedExcitation$ takes the same value on the two members of each pair, and
 \begin{equation}\label{eq.contrastFromSignedExcitation}
  \intensityContrast(t)
  = \sum_{\eone} \pressure_{\eone} \signedExcitation_{\eone}
  \decayedCounts_{\eone}(t) .
 \end{equation}
\end{prop}
\begin{proof}
 Write $u := \excitation^{\transpose}\pressure$.
 Then $\swapMatrix u
 = \big( \swapMatrix \excitation^{\transpose} \swapMatrix \big) \swapMatrix \pressure
 = -u$ by \eqref{eq.directionSymmetry},
 so $u_{\swapMatrix(\eone)} = -u_{\eone}$;
 and $\pressure_{\swapMatrix(\eone)} = -\pressure_{\eone}$,
 so the ratio \eqref{eq.signedExcitation} is unchanged by the swap.
 Equation \eqref{eq.contrastFromSignedExcitation} is
 \eqref{eq.contrastFromState} written entry by entry.
\end{proof}

We read $\signedExcitation_{\eone}$ as the signed excitation strength of the type $\eone$:
one event of type $\eone$ raises the intensity of the types that push its own way, relative
to those that push the other way, by $\signedExcitation_{\eone}$.
A type with $\signedExcitation_{\eone} > 0$ begets its own pressure,
one with $\signedExcitation_{\eone} < 0$ begets the opposite,
and one with $\signedExcitation_{\eone} = 0$ begets neither.
The sign is a property of the kernel and not of the type:
the same six types carry either sign under different specifications.

The intensity contrast and the order flow imbalance are statistics of the same past events,
weighted differently.

By \eqref{eq.contrastFromState} the contrast is a function of the state $\decayedCounts$ and
of the excitation matrix.
It is available only to an observer who knows $\excitation$ and $\decay$,
and under direction symmetry it is blind to $\baseIntensity$,
so it says nothing about the immigrants that are about to arrive.
It reads the types of the past events and weights them by how long ago they occurred.

The imbalance is model free.
By Proposition \ref{prop.signedSizeContribution} it reads the same past events,
but weights them by their sizes and by nothing else:
every event in the window counts its own $\orderSize_n$, whatever its age,
and every event outside the window counts zero.
The kernel governs when events happen and of which type;
the marks govern how far they move the price.
No function of $\decayedCounts$ recovers a size.
